\sectionnn{Conclusion et Perspectives}

Ce travail de stage a permis de concevoir, d'implémenter et de comparer différentes approches de résolution exactes pour le problème d'ordonnancement de projet sous contraintes de ressources multi-compétences avec flexibilité d'affectation. Nous avons formulé un modèle mathématique initial en PLNE, puis deux modèles monolithiques en programmation par contraintes exploitant des variables d'intervalles unitaires pour capturer la flexibilité d'affectation dynamique. Face aux limites de ces approches  (explosion du nombre de variables optionnelles et affaiblissement de la propagation), nous avons développé un algorithme de décomposition de Benders basé sur la logique couplant un problème maître en CP et des sous-problèmes d'affectation résolus sous forme de flot maximum. Les résultats expérimentaux sur les instances de référence confirment l'efficacité de la décomposition LBBD, qui surpasse les modèles monolithiques.
